\documentclass[main]{subfiles}

\begin{document}

\chapter{Constraint Database Compilation and Matter--Antimatter Comparison}

\section{Overview}

Chapters~2 and~3 established the theoretical machinery connecting SME
coefficients to Dobrescu--Mocioiu (DM) potentials. This chapter applies that
machinery to real experimental data. Section~\ref{sec:methodology} describes
the compilation methodology implemented in the SPINDEP framework built for
this thesis: dataset discovery and classification, unit standardisation,
and matter--antimatter pairing. Section~\ref{sec:database} summarises the
resulting database of 273 experimental constraint datasets. Section~\ref{sec:asymmetry}
presents the ten matter--antimatter pairs SPINDEP was able to match, the
full statistical analysis of each (including an autocorrelation-corrected
effective degrees of freedom and a bootstrap confidence interval on the
asymmetry parameter), and their interpretation under the sensitivity-gap
caveat established in Chapter~3.

\section{Database Compilation Methodology}
\label{sec:methodology}

\subsection{Dataset Discovery and Classification}

Each experimental constraint is stored as a two-column CSV file (interaction
range $\lambda$ in metres, upper bound on $|g|$) inside a directory tree
organised by coupling family and interaction class, following the naming
convention
\begin{center}
\texttt{\{V\}\{Author\}\{Year\}\_m\_abs\_\{sector\}.csv}
\end{center}
A parser recursively discovers every CSV under the dataset root, extracts
the potential label, source citation, and fermion sector from the filename
and directory path, and resolves alternate spellings of antimatter sectors
(for example \texttt{ebare}, \texttt{ebar}, and \texttt{epbare} are all
resolved to the canonical labels \texttt{eebar} and \texttt{epbar}). Files
whose sector token cannot be resolved are still retained, with their
potential recorded as \texttt{UNKNOWN} rather than being silently dropped,
so that the size of this classification gap remains visible in the
database rather than hidden by omission.

\subsection{Unit Standardisation}

Interaction ranges are reported across the literature in a mix of units
(metres, MeV$^{-1}$, eV$^{-1}$, centimetres, nanometres, femtometres). A
unit-detection routine inspects each filename for an explicit unit token
and applies the corresponding conversion factor to metres before any
dataset is used in downstream analysis, so that every constraint curve in
the database is directly comparable on a common $\lambda$ axis regardless
of the units used in its original publication.

\subsection{Matter--Antimatter Pairing}

Two datasets are matched into a matter--antimatter pair when they share the
same coupling family, the same DM potential label, and the same interaction
class, and their fermion sectors are physically conjugate (for example
\texttt{ee} with \texttt{eebar}), with exactly one of the two flagged as an
antimatter dataset. This pairing logic is deliberately conservative: it
will not match two datasets that differ in potential classification even
if their sectors are conjugate, since doing so would conflate constraints
on physically distinct interactions.

\section{The Compiled Database}
\label{sec:database}

\subsection{Summary Statistics}

Table~\ref{tab:coupling-breakdown} summarises the 273 datasets compiled for
this thesis by coupling family. Two entries are classified under a coupling
label of \texttt{Code-plot-matlab}, which is a parser artifact rather than a
genuine physical coupling family, arising from a source directory named for
its file format rather than its physics content; these two entries are
excluded from the physics analysis in this chapter and flagged in
Chapter~5 as a data-hygiene item for future cleanup, rather than silently
corrected or removed from the reported count.

\begin{table}[h]
\centering
\begin{tabular}{@{}lr@{}}
\toprule
Coupling family & Datasets \\
\midrule
$g_Ag_A$ & 63 \\
$g_pg_s$ & 53 \\
$g_sg_s$ & 49 \\
$g_Vg_V$ & 43 \\
$g_Ag_V$ & 37 \\
$g_pg_p$ & 15 \\
$V_1$ (unlabelled coupling) & 11 \\
\texttt{Code-plot-matlab} (parser artifact) & 2 \\
\midrule
Total & 273 \\
\bottomrule
\end{tabular}
\caption{Compiled dataset count by coupling family.}
\label{tab:coupling-breakdown}
\end{table}

By DM potential, the database is dominated by $V_{4+5}$ (61 datasets),
followed by potentials that the parser could not resolve to a specific
$V_n$ token (50 datasets, recorded as \texttt{UNKNOWN}), $V_{9+10}$ (31),
$V_{2+3}$ (30), $V_{1a}$ (29), $V_2$ (20), $V_1$ (19, including two entries
reclassified from \texttt{UNKNOWN}; see below), and $V_{12+13}$ (16), with
the remaining potentials each represented by fewer than ten datasets. The
50 remaining \texttt{UNKNOWN}-potential datasets are a genuine
classification gap discussed further in Chapter~5, not an artifact of this
thesis's own tooling: their filenames simply do not contain a token the
naming convention can resolve. Two datasets that previously fell into this
category -- the \texttt{Delaunay\_2017} and \texttt{Adkins\_2022\_eeplus}
files used in Section~\ref{sec:asymmetry} below -- have since been
confirmed as $V_1$ (monopole--monopole, spin-independent scalar) by
checking their source publications directly: Delaunay, Frugiuele, Fuchs and
Soreq (2017) constrain a spin-independent scalar interaction between
electrons, and Adkins, Cassidy and P\'{e}rez-R\'{i}os (2022) report a bound
on the analogous spin-independent $g_s^eg_s^{e^+}$ coupling from
positronium spectroscopy -- both a direct match to $V_1$'s definition, not
an inference from the $g_sg_s$ coupling label alone.

\subsection{Antimatter-Sector Coverage}

Of the 273 datasets, only 17 involve an antimatter sector: 9 in the
$e$-$\bar p$ (\texttt{epbar}) sector, 5 in $e$-$e^+$ (\texttt{eebar}), and 3
in $e$-$\bar\mu$ (\texttt{emubar}). No dataset in the compiled database
touches the $n$-$\bar n$, $n$-$\bar p$, $\mu$-$\bar\mu$, $p$-$\bar p$, or
$e$-$\bar N$ sectors at all. This stark asymmetry -- 17 antimatter datasets
against 256 matter-sector datasets -- is itself the central empirical fact
this thesis works with: it is the reason a meaningful CPT comparison is
currently possible in only a handful of channels, and it directly motivates
the gap analysis of Chapter~5.

\section{Matter--Antimatter Asymmetry Analysis}
\label{sec:asymmetry}

\subsection{Statistical Methodology}

For each matched pair, both constraint curves are interpolated (log-linear
in $\log\lambda$--$\log g$) onto a common log-spaced grid of 300 points
spanning their overlapping $\lambda$ range, and the asymmetry parameter
\begin{equation}
  A_\alpha(\lambda) = \frac{g_m(\lambda) - g_{\bar a}(\lambda)}{g_m(\lambda) + g_{\bar a}(\lambda)}
\end{equation}
is evaluated at every grid point. Statistical significance is assessed four
ways, each addressing a limitation of the one before it. A uniform
$\chi^2$ test assumes a flat 10\% fractional uncertainty at every point. A
weighted $\chi^2$ test instead estimates a per-point fractional uncertainty
from the local curvature of each constraint curve in log-log space, on the
grounds that a smoothly-varying published bound is better determined than
one with sharp local features. Both of these treat the 300 grid points as
independent degrees of freedom, which is not justified: the points are
interpolated from a much smaller number of real measurements along smooth
curves, so adjacent points are strongly correlated. An effective degrees of
freedom is therefore estimated from the autocorrelation length of the
$\chi^2$ residuals in grid-index space (the lag at which the residual
autocorrelation first drops below $1/e$), and the weighted $\chi^2$
significance is recomputed against this corrected, typically much smaller,
degrees of freedom. Finally, a parametric bootstrap (2000 resamples,
propagating the estimated per-point uncertainty through to $A_\alpha$
itself) gives a 95\% confidence interval on the mean $|A_\alpha|$ for each
pair.

\subsection{The Ten Matched Pairs}

Table~\ref{tab:pairs} gives the full results for every matter--antimatter
pair SPINDEP was able to match from the compiled database.

\begin{table}[h]
\centering
\small
\begin{tabular}{@{}llllrrrr@{}}
\toprule
Coupling & Potential & Sector & Matter / antimatter source & $\overline{|A_\alpha|}$ & dof$_{\text{eff}}$ & $p_{\text{eff}}$ & 95\% CI \\
\midrule
$g_Ag_A$    & $V_2$    & $e$-$\bar p$ & Karshenboim2011 / Ficek2018     & 0.9998 & 17 & $\approx0$ & [0.9994, 0.9996] \\
$g_Ag_A$    & $V_2$    & $e$-$e^+$    & Ficek2017 / Karshenboim2011     & 0.9892 & 21 & $\approx0$ & [0.9891, 0.9893] \\
$g_Ag_A$    & $V_{2+3}$& $e$-$e^+$    & Ficek2017 / Fadeev2022          & 0.9539 & 6  & $\approx0$ & [0.9538, 0.9540] \\
$g_pg_p$    & $V_{2+3}$& $e$-$e^+$    & Fadeev2022 / Fadeev2022         & 0.9535 & 6  & $\approx0$ & [0.9534, 0.9537] \\
$g_Vg_V$    & $V_{2+3}$& $e$-$e^+$    & Fadeev2022 / Fadeev2022         & 0.9535 & 6  & $\approx0$ & [0.9534, 0.9537] \\
$g_sg_s$    & $V_1$     & $e$-$e^+$    & Delaunay2017 / Adkins2022       & 0.8727 & 21 & $\approx0$ & [0.8726, 0.8728] \\
$g_Ag_A$    & $V_{2+3}$& $e$-$\bar p$ & Fadeev2022 / Ficek2018          & 0.8237 & 6  & $\approx0$ & [0.8235, 0.8237] \\
$g_Ag_A$    & $V_{2+3}$& $e$-$\bar p$ & Fadeev2022 / Fadeev2022         & 0.8044 & 8  & $\approx0$ & [0.8043, 0.8045] \\
$g_Vg_V$    & $V_{2+3}$& $e$-$\bar p$ & Fadeev2022 / Ficek2018          & 0.7994 & 7  & $\approx0$ & [0.7993, 0.7995] \\
$g_Ag_A$    & $V_2$    & $e$-$e^+$    & Jiao2019 / Karshenboim2011      & 0.3336 & 9  & $\approx0$ & [0.3335, 0.3337] \\
\bottomrule
\end{tabular}
\caption{All ten matched matter--antimatter pairs, with the autocorrelation-corrected effective degrees of freedom, the resulting weighted-$\chi^2$ significance, and the 95\% bootstrap confidence interval on the mean asymmetry.}
\label{tab:pairs}
\end{table}

Effective degrees of freedom range from 6 to 21 across the ten pairs, a
factor of 14--50 below the naive 300 grid points, confirming that the
autocorrelation correction described in Section~\ref{sec:methodology} is not
a negligible refinement: an analysis quoting significance at $\mathrm{dof}=300$
throughout would be substantially overstating its own precision. Every pair
nonetheless remains significant at effectively $p\to0$ even after this
correction, since the $\chi^2$ values themselves (tens of thousands, in the
underlying per-pair data) are large enough relative to even the reduced
degrees of freedom that the correction changes the reported $p$-value by
many orders of magnitude without changing its qualitative conclusion.

\subsection{Interpretation}

Per the argument developed fully in Chapter~3, Section~3.2.3, an
asymmetry value $|A_\alpha|$ near 1 in Table~\ref{tab:pairs} is
\emph{consistent with}, but not \emph{evidence for}, CPT violation in the
corresponding channel: SPINDEP compares independent, one-sided experimental
upper bounds, and a large sensitivity gap between the matter- and
antimatter-sector measurement being compared produces exactly the same
signature regardless of the true underlying physics. This applies uniformly
to every row of Table~\ref{tab:pairs}, including the $g_Ag_A\cdot V_2\cdot e$-$\bar p$
pair at $0.9998$, the highest asymmetry in the database, and the
$g_sg_s\cdot V_1\cdot e$-$e^+$ pair at $0.8727$, whose scalar-scalar
coupling is CPT-even at leading order and whose asymmetry is therefore
attributable entirely to sensitivity gap rather than any sign-flip
mechanism -- a useful internal check, since applying the \emph{same}
reasoning uniformly to the CPT-odd $g_Ag_A$ rows means their high asymmetry
values cannot, by this data alone, be distinguished from the same
effect. Distinguishing a genuine CPT-violating signal from a sensitivity
gap requires either bounds of demonstrably comparable precision on both
sides of a pair, or actual signed central-value measurements rather than
one-sided limits -- neither of which the currently compiled database
provides for any pair. This is the central limitation Chapter~5 addresses
directly.

The lowest asymmetry in the table, $0.334$ for the $g_Ag_A\cdot V_2\cdot e$-$e^+$
pair (Jiao2019 vs.\ Karshenboim2011), is the pair whose matter- and
antimatter-sector bounds happen to be closest in sensitivity among the ten;
it is the single most informative data point currently in the database for
constraining $b_\mu$ in the $e$-$e^+$ channel precisely because it is
\emph{not} saturated at $|A_\alpha|\approx1$.

\subsection{Figures}

Figure~\ref{fig:atlas} shows the combined constraint atlas across the
eleven classified potential labels in the compiled database. These map onto
13 of the 16 catalogue potentials from Chapter~2 -- several labels such as
$V_{2+3}$, $V_{4+5}$, $V_{9+10}$, and $V_{12+13}$ each combine two adjacent
catalogue potentials that the underlying experiments do not distinguish
individually. $V_6$, $V_7$, and $V_{14}$ have no classified dataset at all,
distinct from the 50 \texttt{UNKNOWN}-potential entries discussed above,
which have data but no resolved label. Figure~\ref{fig:matter-antimatter}
shows an example matter--antimatter comparison plot for the
highest-asymmetry pair in Table~\ref{tab:pairs}. The complete set of ten
such comparison plots, one per matched pair, is included in the project
repository under \texttt{figures/matter\_antimatter/}.

\begin{figure}[h]
\centering
\includegraphics[width=0.85\textwidth]{constraint_atlas_all.png}
\caption{Combined constraint atlas across all compiled potentials.}
\label{fig:atlas}
\end{figure}

\begin{figure}[h]
\centering
\includegraphics[width=0.75\textwidth]{gAgA_V2_ep_Karshenboim2011_vs_Ficek2018.png}
\caption{Matter--antimatter comparison for the $g_Ag_A\cdot V_2\cdot e$-$\bar p$ pair (Karshenboim2011 vs.\ Ficek2018), the highest asymmetry in the compiled database.}
\label{fig:matter-antimatter}
\end{figure}

\section{Chapter Summary}

This chapter compiled 273 experimental constraint datasets into a
standardised database, of which only 17 involve an antimatter sector,
concentrated in just three channels ($e$-$\bar p$, $e$-$e^+$, $e$-$\bar\mu$).
Ten matter--antimatter pairs could be matched from this database and
analysed with a statistical methodology that corrects for the correlated
nature of interpolated constraint points. Every pair shows a statistically
significant asymmetry even after this correction, but -- applying the
interpretive framework of Chapter~3 uniformly -- none of these results can
currently be read as evidence of CPT violation rather than an ordinary
sensitivity gap between the matter- and antimatter-sector experiments being
compared. Chapter~5 quantifies exactly where the antimatter-sector coverage
gaps responsible for this limitation are largest, and proposes where future
experimental effort would do the most to close them.

\end{document}
